\documentclass[12pt]{article}

% 必要的包
\usepackage[margin=1in]{geometry}
\usepackage{setspace}
\usepackage{enumitem}
\usepackage{times}
\usepackage{fancyhdr}
\usepackage{graphicx}
\usepackage{hyperref}

% 页码设置
\pagestyle{fancy}
\fancyhf{}
\renewcommand{\headrulewidth}{0pt}
\fancyfoot[C]{\thepage}

% 行间距 - 法院通常要求双倍行距
\doublespacing

% 格式设置
\setlength{\parindent}{0.5in}
\setlength{\parskip}{0pt}

\begin{document}

%------------------------------
% 免责声明 - 提交前请删除
%------------------------------
\begin{center}
\textbf{\large IMPORTANT DISCLAIMER}
\end{center}

\begin{singlespace}
\small
THIS DOCUMENT IS A GENERAL TEMPLATE FOR EDUCATIONAL PURPOSES ONLY. 
IT DOES NOT CONSTITUTE LEGAL ADVICE. WRITS OF MANDAMUS ARE COMPLEX 
LEGAL DOCUMENTS THAT REQUIRE COMPLIANCE WITH SPECIFIC COURT RULES, 
STANDING REQUIREMENTS, AND PROCEDURAL REQUIREMENTS. FILING AN 
IMPROPERLY DRAFTED WRIT MAY RESULT IN YOUR CASE BEING DISMISSED 
OR OTHER ADVERSE CONSEQUENCES. YOU MUST CONSULT A LICENSED 
ATTORNEY IN YOUR JURISDICTION BEFORE FILING ANY LEGAL DOCUMENT 
WITH A COURT. NOTHING IN THIS TEMPLATE CREATES AN ATTORNEY-CLIENT 
RELATIONSHIP.
\end{singlespace}

\vspace{1in}
\newpage

%------------------------------
% 标题页 - 巴尔的摩县巡回法院版本
%------------------------------
\begin{center}
\textbf{\Large IN THE CIRCUIT COURT FOR BALTIMORE COUNTY, MARYLAND}
\vspace{0.5in}

CAUSTIN LEE McLAUGHLIN, and \\
JUAN JOSE FERNANDEZ, \\
\textit{Petitioners,}

\vspace{0.3in}
v.
\vspace{0.3in}

[RESPONDENT NAME 1], in his/her official capacity; \\
[RESPONDENT NAME 2], in his/her official capacity; \\
[ADD ADDITIONAL RESPONDENTS AS NEEDED],
\textit{Respondents.}

\vspace{0.5in}

\textbf{\large PETITION FOR WRIT OF MANDAMUS}
\vspace{0.25in}

Case No.: \underline{\hspace{2in}}
\end{center}

\vspace{0.75in}

%------------------------------
% 签名栏
%------------------------------
\begin{flushleft}
Respectfully submitted,
\vspace{0.5in}

Caustin Lee McLaughlin \\
Juan Jose Fernandez \\
[YOUR ADDRESS] \\
[CITY, STATE ZIP] \\
[TELEPHONE NUMBER] \\
[EMAIL ADDRESS] \\
\textit{Pro Se Petitioners}
\vspace{0.5in}

Date: \underline{\hspace{2in}}
\end{flushleft}

\newpage

%------------------------------
% 证书送达页
%------------------------------
\section*{CERTIFICATE OF SERVICE}

I hereby certify that on \underline{\hspace{1in}}, a true and correct copy of the foregoing Petition for Writ of Mandamus was served on the following parties in the manner indicated:

\begin{enumerate}[label=\arabic*.]
    \item [NAME], [TITLE/ADDRESS] - [Method of service]
    \item [NAME], [TITLE/ADDRESS] - [Method of service]
\end{enumerate}

\vspace{0.5in}

\underline{\hspace{3in}} \\
Caustin Lee McLaughlin

\vspace{0.25in}

Date: \underline{\hspace{2in}}

\newpage

%------------------------------
% 目录
%------------------------------
\tableofcontents
\newpage

%------------------------------
% 管辖权陈述
%------------------------------
\section{STATEMENT OF JURISDICTION}

Petitioners Caustin Lee McLaughlin and Juan Jose Fernandez file this Petition for Writ of Mandamus pursuant to Maryland Rule 15-701 and the common law power of this Court to issue writs of mandamus to compel the performance of ministerial acts by public officials. This Court has jurisdiction over this matter because:

\begin{enumerate}[label=(\alph*)]
    \item Respondents are public officials/entities within Baltimore County, Maryland;
    \item The acts sought to be compelled are ministerial duties owed to Petitioners;
    \item Petitioners have a clear legal right to the performance of the duties sought;
    \item Petitioners have no plain, speedy, and adequate remedy at law; and
    \item Petitioners have given Respondents proper notice of this Petition.
\end{enumerate}

%------------------------------
% 当事人陈述
%------------------------------
\section{PARTIES}

\begin{enumerate}[label=(\alph*)]
    \item \textbf{Petitioner Caustin Lee McLaughlin} is an individual residing at [ADDRESS], Baltimore County, Maryland. Petitioner McLaughlin seeks Schedule A/Schedule C employment as described herein as part of the implementation referenced in supporting materials.
    
    \item \textbf{Petitioner Juan Jose Fernandez} is an individual residing at [ADDRESS], Baltimore County, Maryland, and is the spouse of Petitioner McLaughlin. Petitioner Fernandez seeks Schedule A/Schedule C employment as described herein.
    
    \item \textbf{Respondent [NAME]} is [TITLE] of [AGENCY], located at [ADDRESS], Baltimore County, Maryland. Respondent is sued in his/her official capacity.
    
    \item [ADD ADDITIONAL RESPONDENTS AS APPROPRIATE]
\end{enumerate}

%------------------------------
% 事实陈述
%------------------------------
\section{STATEMENT OF FACTS}

[INSERT DETAILED STATEMENT OF FACTS HERE. Include the following:
\begin{enumerate}
    \item Background of Petitioners' eligibility for Schedule A/Schedule C hiring
    \item The specific ministerial duties Respondents have failed to perform
    \item Timeline of requests made by Petitioners
    \item Any administrative remedies exhausted
    \item How Petitioners have been harmed by the failure to act
    \item Reference to the larger implementation program referenced in your materials
\end{enumerate}]

%------------------------------
% 理由
%------------------------------
\section{ARGUMENT}

\subsection{Standard for Issuance of Writ of Mandamus in Maryland}

Under Maryland law, a writ of mandamus is appropriate to compel the performance of a ministerial act where: (1) the petitioner has a clear legal right to the relief sought; (2) the respondent has a clear ministerial duty to perform the act; and (3) there is no other plain, speedy, and adequate remedy at law. \textit{See, e.g.,}, \textit{Heck v. Chesapeake \& Potomac Tel. Co.}, 288 Md. 523 (1980).

\subsection{Petitioners Have a Clear Legal Right to Schedule A/Schedule C Hiring Consideration/Appointment}

[INSERT LEGAL ARGUMENT REGARDING SCHEDULE A/SCHEDULE C ELIGIBILITY AND ENTITLEMENT]

\subsection{Respondents Have a Clear Ministerial Duty to Act}

[INSERT ARGUMENT THAT THE DUTY IS MINISTERIAL, NOT DISCRETIONARY]

\subsection{Petitioners Have No Adequate Remedy at Law}

[EXPLAIN WHY OTHER REMEDIES (E.G., ADMINISTRATIVE APPEALS, DAMAGE CLAIMS) ARE INADEQUATE]

%------------------------------
% 请求救济
%------------------------------
\section{PRAYER FOR RELIEF}

WHEREFORE, Petitioners Caustin Lee McLaughlin and Juan Jose Fernandez respectfully request that this Court:

\begin{enumerate}[label=\arabic*.]
    \item Issue a Writ of Mandamus directing Respondents to [SPECIFIC ACTS SOUGHT, e.g., "process Petitioners' Schedule A hiring applications in accordance with applicable law and regulations"];
    
    \item Order that Petitioners Caustin Lee McLaughlin and Juan Jose Fernandez be considered/appointed to appropriate positions pursuant to Schedule A/Schedule C hiring authorities as part of the referenced implementation program;
    
    \item Grant such other and further relief as the Court deems just and proper.
\end{enumerate}

\vspace{0.5in}

Respectfully submitted,

\vspace{0.5in}

\underline{\hspace{3in}} \\
Caustin Lee McLaughlin \\
Pro Se Petitioner

\vspace{0.3in}

\underline{\hspace{3in}} \\
Juan Jose Fernandez \\
Pro Se Petitioner

\newpage

%------------------------------
% 附录 - 注意：这是示例结构
%------------------------------
\section*{EXHIBIT A - SCHEDULE A/SCHEDULE C POSITION REQUESTS}
[DESCRIBE POSITIONS SOUGHT FOR CAUSTIN LEE McLAUGHLIN AND JUAN JOSE FERNANDEZ AS PART OF LARGER IMPLEMENTATION]

\section*{EXHIBIT B - SUPPORTING DOCUMENTATION}
[LIST ATTACHED DOCUMENTS - your supporting materials, correspondence, eligibility documentation, etc.]

\newpage

%==============================================================
% DC CIRCUIT COURT 版本备注
%==============================================================
\begin{center}
\textbf{\Large NOTE REGARDING D.C. CIRCUIT COURT FILING}
\end{center}

For petitions to the United States Court of Appeals for the District of Columbia Circuit (for writs of mandamus against federal officials/agencies), note the following differences:

\begin{enumerate}
    \item \textbf{Authority}: Fed. R. App. P. 21; D.C. Cir. Rule 21
    \item \textbf{Caption}: Should read "IN THE UNITED STATES COURT OF APPEALS FOR THE DISTRICT OF COLUMBIA CIRCUIT"
    \item \textbf{Parties}: If against federal officials, identify proper federal respondents (e.g., agency heads in official capacity)
    \item \textbf{Venue}: Ensure D.C. Circuit has proper jurisdiction over the respondents/actions
    \item \textbf{Format}: Follow Federal Rules of Appellate Procedure and D.C. Circuit Local Rules regarding page limits, formatting, and filing requirements
    \item \textbf{Filing fee}: Check current filing fee requirements
\end{enumerate}

\vspace{0.5in}

\textbf{CRITICAL REMINDER}: This template is for educational purposes only. Schedule A and Schedule C hiring authorities are federal personnel mechanisms, and mandamus relief in federal court involves complex questions of sovereign immunity, standing, and administrative law. You MUST consult an attorney admitted to practice in the relevant courts before filing anything.

\end{document}
